\chapter{Системы сборки, базы данных и функциональный стиль}
\label{ch:06}

До сих пор программы из этой книги жили сами по себе: мы компилировали один
файл вручную и держали данные в памяти. Реальный проект устроен иначе~--- в
нём десятки сторонних библиотек, набор тестов и база данных, переживающая
перезапуск приложения. Эта глава закрывает разрыв между учебным примером и
рабочим проектом: мы поручим сборку системам Maven и Gradle, научимся
разговаривать с базой данных на языке JDBC и на языке объектов через
Hibernate, разберёмся с транзакциями. Завершает главу Stream~API~---
декларативный способ обработать данные, не написав ни одного цикла; коллекции,
на которые он опирается, разбирались в главе~\ref{ch:05}.

\section{Системы автоматической сборки}

\subsection{Maven: структура проекта и файл pom.xml}

Кроме написания кода с проектом приходится делать многое: скачивать JAR-файлы
библиотек и всех библиотек, от которых зависят они сами; прописывать пути к
ним при вызове компилятора; запускать тесты; упаковывать результат в архив.
Вручную это выполнимо ровно один раз~--- на второй день превращается в пытку.
\term{Система сборки} (\eng{build system})~--- автоматизированная кухня, где
ингредиенты доставляются сами. Главная идея самой распространённой из них,
Maven,~--- \term{соглашение важнее настройки} (\eng{convention over
configuration}): инструмент навязывает проекту единую раскладку папок
(рис.~\ref{fig:06-1}) и за это сам знает, где лежат исходники, тесты и
ресурсы.

\begin{figure}[htbp]\centering
\begin{forest} hierarchy
[project/
  [pom.xml]
  [src/
    [main/
      [java/]
      [resources/]
    ]
    [test/
      [java/]
      [resources/]
    ]
  ]
  [target/
    [classes/]
    [project-1.0.jar]
  ]
]
\end{forest}
\caption{Стандартная структура Maven-проекта}\label{fig:06-1}
\end{figure}

Каталог \texttt{target} создаётся при сборке и в систему контроля версий не
отправляется: он целиком восстанавливается из исходников. Сердце проекта~---
файл \texttt{pom.xml} (\eng{Project Object Model}), отвечающий на три вопроса:
чем проект является, из чего собран и как его собирать
(листинг~\ref{lst:06-1}). Тройка \texttt{groupId}, \texttt{artifactId} и
\texttt{version}~--- это \term{GAV-координаты}, уникальный адрес артефакта.

\begin{listing}[H]
\begin{minted}{xml}
<project xmlns="http://maven.apache.org/POM/4.0.0">
  <modelVersion>4.0.0</modelVersion>

  <!-- GAV-координаты: группа, артефакт, версия -->
  <groupId>com.example</groupId>
  <artifactId>movie-app</artifactId>
  <version>1.0-SNAPSHOT</version>

  <properties>
    <maven.compiler.source>21</maven.compiler.source>
    <maven.compiler.target>21</maven.compiler.target>
  </properties>

  <dependencies>
    <dependency>
      <groupId>com.h2database</groupId>
      <artifactId>h2</artifactId>
      <version>2.2.224</version>
    </dependency>
    <dependency>
      <groupId>org.junit.jupiter</groupId>
      <artifactId>junit-jupiter</artifactId>
      <version>5.10.0</version>
      <scope>test</scope>   <!-- нужна только тестам -->
    </dependency>
  </dependencies>
</project>
\end{minted}
\caption{Файл pom.xml: координаты, свойства и зависимости}
\label{lst:06-1}
\end{listing}

Элемент \texttt{scope} отвечает на вопрос «на каком этапе жизни проекта эта
библиотека нужна»: \texttt{compile} (по умолчанию)~--- везде;
\texttt{test}~--- только при компиляции и запуске тестов, в готовый JAR такая
библиотека не попадёт; \texttt{provided}~--- при компиляции, а в среде
выполнения её даст сервер приложений; \texttt{runtime}~--- только при запуске
(типичный случай~--- драйвер СУБД). Зависимости Maven ищет сначала в локальном
кэше (папка \texttt{.m2/repository}), затем в репозитории Maven~Central,
поэтому вторая сборка идёт без сети.

\subsection{Жизненный цикл Maven}

Сборка разбита на \term{фазы} (\eng{phases}), идущие строго друг за другом
(табл.~\ref{tab:06-1}). Правило одно, и оно объясняет всё поведение
инструмента: запуская фазу, вы автоматически запускаете все предшествующие.
Поэтому \texttt{mvn package} не просто складывает классы в архив: она проверит
проект, скомпилирует его и прогонит тесты, и если хоть один тест упадёт,
JAR-файла вы не получите.

\begin{table}[htbp]
\caption{Основные фазы жизненного цикла Maven}
\label{tab:06-1}
\centering
\begin{tabular}{L{2.6cm}L{12cm}}
\toprule
Фаза & Что происходит \\
\midrule
\texttt{validate} & проверка корректности проекта \\
\texttt{compile} & компиляция исходного кода в \texttt{target/classes} \\
\texttt{test} & компиляция и запуск модульных тестов \\
\texttt{package} & упаковка результата в JAR или WAR \\
\texttt{install} & установка артефакта в локальный репозиторий \\
\texttt{deploy} & публикация артефакта в удалённый репозиторий \\
\bottomrule
\end{tabular}
\end{table}

Команд, которых хватает на весь курс, шесть (листинг~\ref{lst:06-2}). Слово
\texttt{clean}~--- не фаза цикла, а отдельная команда, удаляющая каталог
\texttt{target}.

\begin{listing}[H]
\begin{minted}{bash}
mvn compile          # скомпилировать исходники
mvn test             # компиляция + запуск тестов
mvn package          # + упаковка в JAR (файл появится в target/)
mvn clean package    # удалить target/, затем собрать заново
mvn install          # + положить артефакт в локальный репозиторий
mvn dependency:tree  # показать дерево зависимостей проекта
\end{minted}
\caption{Команды Maven, которые нужны каждый день}
\label{lst:06-2}
\end{listing}

Последняя команда решает проблему, о которой начинающие не подозревают.
Подключённая библиотека тянет за собой собственные зависимости, те~--- свои;
такие зависимости называются \term{транзитивными}. Одна строчка в
\texttt{pom.xml} легко превращается в сорок JAR-файлов, а две библиотеки могут
потребовать разных версий одной и той же третьей.

\subsection{Gradle}

Если \texttt{pom.xml}~--- подробная декларация, где каждый шаг расписан
по-XML-евски многословно, то файл сборки Gradle~--- лаконичный исполняемый
скрипт на языке Groovy или Kotlin (листинг~\ref{lst:06-3}). Раскладка
каталогов та же, что у Maven.

\begin{listing}[H]
\begin{minted}{kotlin}
plugins { java }

repositories {
    mavenCentral()          // откуда брать зависимости
}

dependencies {
    // implementation - аналог scope compile у Maven
    implementation("com.h2database:h2:2.2.224")
    // testImplementation - аналог scope test
    testImplementation("org.junit.jupiter:junit-jupiter:5.10.0")
}
\end{minted}
\caption{Файл build.gradle.kts}
\label{lst:06-3}
\end{listing}

Рядом с ним в корне лежат \texttt{settings.gradle.kts} (состав модулей) и пара
\texttt{gradlew}, \texttt{gradlew.bat}~--- \term{обёртка Gradle}
(\eng{Gradle Wrapper}), которая сама скачивает записанную в проекте версию
Gradle и запускает сборку ею: устанавливать Gradle участникам команды не
нужно. Выбор между двумя системами практический: Maven предсказуем и потому
остаётся стандартом в корпоративной разработке; Gradle гибче и быстрее на
больших проектах за счёт инкрементальной сборки, а под Android обязателен.

\section{Работа с базой данных: JDBC}

\subsection{Архитектура JDBC}

\term{JDBC} (\eng{Java Database Connectivity})~--- стандартный интерфейс Java
для работы с реляционными базами данных. Ключевое слово здесь
«интерфейс»: пакет \texttt{java.sql} не содержит кода, умеющего разговаривать
с конкретной СУБД, он лишь описывает, какие методы у такого кода должны быть,
а реализацию~--- \term{драйвер}~--- поставляет производитель базы. Представьте
розетку: стандарт один, а электростанции у всех разные. Цепочка от приложения
до базы показана на рис.~\ref{fig:06-2}; при переезде с H2 на PostgreSQL в ней
меняются ровно два звена~--- драйвер и строка подключения.

\begin{figure}[htbp]\centering
\begin{tikzpicture}[font=\small,node distance=5mm]
\node[boxw=68mm] (a) {Java-приложение};
\node[boxw=68mm,below=of a] (b) {JDBC API~--- интерфейсы пакета \texttt{java.sql}};
\node[boxw=68mm,below=of b] (c) {\texttt{DriverManager}~--- выбирает драйвер по URL};
\node[boxw=68mm,below=of c] (d) {Драйвер JDBC: H2, PostgreSQL, MySQL};
\node[boxw=68mm,below=of d] (e) {База данных};
\draw[flowarrow] (a)--(b);\draw[flowarrow] (b)--(c);
\draw[flowarrow] (c)--(d);\draw[flowarrow] (d)--(e);
\end{tikzpicture}
\caption{Архитектура доступа к базе данных через JDBC}\label{fig:06-2}
\end{figure}

Интерфейсов, которыми пользуются каждый день, пять (табл.~\ref{tab:06-2}). В
примерах ниже используется база H2: она написана на Java, подключается одной
зависимостью и работает в оперативной памяти, не требуя установки сервера.

\begin{table}[htbp]
\caption{Основные интерфейсы JDBC}
\label{tab:06-2}
\centering
\begin{tabular}{L{4.2cm}L{10.4cm}}
\toprule
Интерфейс & Назначение \\
\midrule
\texttt{DriverManager} & выдаёт соединение по строке подключения \\
\texttt{Connection} & соединение с базой; управляет транзакцией \\
\texttt{Statement} & выполняет запрос без параметров \\
\texttt{PreparedStatement} & предкомпилированный запрос с параметрами \\
\texttt{ResultSet} & курсор по строкам результата \texttt{SELECT} \\
\bottomrule
\end{tabular}
\end{table}

\subsection{Полный цикл операций}

Работа с базой идёт по одной схеме: получить соединение, подготовить запрос,
выполнить его, разобрать результат, закрыть всё в обратном порядке. Последний
шаг берёт на себя \texttt{try}-с-ресурсами: все три интерфейса реализуют
\texttt{AutoCloseable}. Знаки вопроса в тексте запросов
(листинг~\ref{lst:06-4})~--- места для параметров.

\begin{listing}[H]
\begin{minted}{java}
// H2 в оперативной памяти: сервер устанавливать не нужно
String url = "jdbc:h2:mem:moviedb;DB_CLOSE_DELAY=-1";

try (Connection conn =
         DriverManager.getConnection(url, "sa", "")) {

    try (Statement st = conn.createStatement()) {
        st.execute("CREATE TABLE movies ("
                 + "id INT AUTO_INCREMENT PRIMARY KEY,"
                 + "title VARCHAR(200) NOT NULL,"
                 + "genre VARCHAR(100), year INT)");
    }
    // Вставка: параметры нумеруются с единицы
    try (PreparedStatement ps = conn.prepareStatement(
            "INSERT INTO movies (title, genre, year)"
          + " VALUES (?, ?, ?)")) {
        ps.setString(1, "Матрица");
        ps.setString(2, "Фантастика");
        ps.setInt(3, 1999);
        ps.executeUpdate();   // вернёт число изменённых строк
    }
    // Выборка: к столбцам обращаемся по имени
    try (PreparedStatement ps = conn.prepareStatement(
            "SELECT * FROM movies WHERE year > ?")) {
        ps.setInt(1, 1990);
        try (ResultSet rs = ps.executeQuery()) {
            while (rs.next()) {
                System.out.println(rs.getString("title"));
            }
        }
    }
}
\end{minted}
\caption{Создание таблицы и операции CRUD через JDBC}
\label{lst:06-4}
\end{listing}

Запомните разделение: \texttt{executeQuery} применяется к \texttt{SELECT} и
возвращает \texttt{ResultSet}; \texttt{executeUpdate}~--- к \texttt{INSERT},
\texttt{UPDATE} и \texttt{DELETE} и возвращает количество затронутых строк.
Обход результата всегда выглядит как \texttt{while (rs.next())}: сразу после
запроса курсор стоит \emph{перед} первой строкой.

\begin{oshibka}
\textbf{Частая ошибка: нумерация с нуля.} И параметры запроса, и столбцы
результата в JDBC нумеруются \emph{с единицы}: вызов \texttt{rs.getString(0)}
бросит \texttt{SQLException}. Обращение по имени столбца надёжнее.
\end{oshibka}

\subsection{SQL-инъекция и параметризованный запрос}

Теперь понятно, зачем нужны знаки вопроса. Соблазн собрать запрос
конкатенацией строк велик, но именно так возникает \term{SQL-инъекция}
(\eng{SQL injection})~--- атака, при которой пользовательский ввод перестаёт
быть данными и становится частью команды (листинг~\ref{lst:06-5}).

\begin{listing}[H]
\begin{minted}{java}
String userInput = "'; DROP TABLE movies; --";

// ОПАСНО: ввод пользователя склеен с текстом запроса
String bad = "SELECT * FROM movies WHERE title = '"
           + userInput + "'";
// В базу уйдёт два запроса, второй удалит таблицу:
// SELECT * FROM movies WHERE title = ''; DROP TABLE movies; --'

// БЕЗОПАСНО: параметр передаётся отдельно от текста запроса
try (PreparedStatement ps = conn.prepareStatement(
        "SELECT * FROM movies WHERE title = ?")) {
    ps.setString(1, userInput);   // это данные, а не код
    ResultSet rs = ps.executeQuery();
}
\end{minted}
\caption{Опасная конкатенация и безопасный параметр}
\label{lst:06-5}
\end{listing}

Разница принципиальна: \texttt{PreparedStatement} отправляет базе шаблон
запроса и значения параметров по отдельности. База разбирает шаблон один раз,
уже зная, где в нём команда, а где данные, и никакая кавычка внутри значения
структуру запроса изменить не может; побочная выгода~--- скорость. Правило:
пользовательский ввод в текст SQL-запроса не подставляется никогда.

\subsection{Паттерн DAO}

В листинге~\ref{lst:06-4} SQL перемешан с логикой программы. На двадцати
строках это незаметно, на десяти тысячах~--- катастрофа: чтобы поменять
структуру таблицы, придётся искать запросы по всему проекту. Ответ~---
\term{DAO} (\eng{Data Access Object}): весь SQL собирается в отдельном слое, а
остальная программа работает с обычными Java-объектами
(листинг~\ref{lst:06-6}).

\begin{listing}[H]
\begin{minted}{java}
public interface MovieDAO {
    void insert(Movie movie) throws SQLException;
    Optional<Movie> findById(int id) throws SQLException;
    List<Movie> findByGenre(String genre) throws SQLException;
    void delete(int id) throws SQLException;
}

public class MovieDAOImpl implements MovieDAO {

    private final Connection conn;

    public MovieDAOImpl(Connection conn) { this.conn = conn; }

    @Override
    public void insert(Movie movie) throws SQLException {
        String sql = "INSERT INTO movies (title, genre, year)"
                   + " VALUES (?, ?, ?)";
        try (PreparedStatement ps = conn.prepareStatement(
                sql, Statement.RETURN_GENERATED_KEYS)) {
            ps.setString(1, movie.getTitle());
            ps.setString(2, movie.getGenre());
            ps.setInt(3, movie.getYear());
            ps.executeUpdate();
            // база сама выдала значение id - забираем его
            try (ResultSet keys = ps.getGeneratedKeys()) {
                if (keys.next()) movie.setId(keys.getInt(1));
            }
        }
    }
}
\end{minted}
\caption{Интерфейс DAO и фрагмент его реализации}
\label{lst:06-6}
\end{listing}

Выгод три: бизнес-логика ничего не знает ни о SQL, ни о JDBC; в тестах
реализацию подменяют заглушкой, и база не нужна; саму реализацию можно
заменить целиком, перейдя на Hibernate. Метод \texttt{findById} возвращает
\texttt{Optional<Movie>}, потому что отсутствие строки~--- обычный результат
поиска, а не ошибка. И ещё правило: превращение строки результата в объект
выносят в единственный приватный метод.

\section{ORM и Hibernate}

\subsection{Сущность и конфигурация}

Даже с паттерном DAO значительная часть кода остаётся однообразной: разложить
объект по параметрам запроса, собрать объект из строки результата, повторить
для каждой сущности. \term{ORM} (\eng{Object-Relational Mapping})~---
технология, которая делает это за вас: вы один раз описываете соответствие
между классом и таблицей, а SQL фреймворк генерирует сам. Самая
распространённая реализация для Java~--- \term{Hibernate}; соответствие
задаётся аннотациями прямо в классе (листинг~\ref{lst:06-7}).

\begin{listing}[H]
\begin{minted}{java}
import jakarta.persistence.*;

@Entity                       // класс отображается на таблицу
@Table(name = "movies")
public class Movie {

    @Id                       // первичный ключ
    @GeneratedValue(strategy = GenerationType.IDENTITY)
    private Integer id;       // значение выдаёт сама база

    @Column(name = "title", nullable = false, length = 200)
    private String title;

    private String genre;
    private Integer year;

    // Конструктор без аргументов обязателен для JPA
    public Movie() {}

    // Геттеры, сеттеры и toString() опущены
}
\end{minted}
\caption{Класс-сущность с аннотациями JPA}
\label{lst:06-7}
\end{listing}

Аннотации отвечают на вопрос «что отображать». На вопрос «куда подключаться и
что делать со схемой» отвечает файл \texttt{hibernate.cfg.xml}
(листинг~\ref{lst:06-8}) из каталога \texttt{src/main/resources}.

\begin{listing}[H]
\begin{minted}{xml}
<hibernate-configuration>
  <session-factory>
    <property name="hibernate.connection.url">
      jdbc:h2:mem:moviedb;DB_CLOSE_DELAY=-1</property>
    <property name="hibernate.connection.username">sa</property>
    <property name="hibernate.dialect">
      org.hibernate.dialect.H2Dialect</property>
    <!-- create: пересоздать схему при каждом запуске;
         update - дополнить, validate - только проверить -->
    <property name="hibernate.hbm2ddl.auto">create</property>
    <property name="hibernate.show_sql">true</property>
    <mapping class="com.movies.Movie"/>
  </session-factory>
</hibernate-configuration>
\end{minted}
\caption{Конфигурация Hibernate для базы H2}
\label{lst:06-8}
\end{listing}

Диалект~--- это знание Hibernate о том, чем SQL конкретной СУБД отличается от
стандарта: как называется автоинкремент, как ограничить число строк. Именно он
позволяет одному коду работать и с H2, и с PostgreSQL. Свойство
\texttt{hbm2ddl.auto} со значением \texttt{create} удобно на занятиях и
недопустимо в рабочей системе: при каждом запуске схема со всеми данными
создаётся заново.

\subsection{SessionFactory, Session и запросы}

\texttt{SessionFactory}~--- тяжёлая фабрика: она читает конфигурацию,
разбирает аннотации сущностей, поднимает пул соединений; создаётся один раз за
время жизни приложения и потокобезопасна. \texttt{Session}~--- лёгкий объект
на одну единицу работы: открыли, сделали дело, закрыли. Аналогия~--- банк и
визит в банк: здание с сейфами строят один раз, а приходят в него ежедневно и
ненадолго. В листинге~\ref{lst:06-9} собраны основные операции; ни одного
запроса \texttt{INSERT} или \texttt{SELECT} мы не пишем.

\begin{listing}[H]
\begin{minted}{java}
SessionFactory factory = new Configuration()
        .configure("hibernate.cfg.xml")
        .buildSessionFactory();

// СОХРАНЕНИЕ: изменения уходят в базу при коммите
try (Session session = factory.openSession()) {
    Transaction tx = session.beginTransaction();
    session.persist(new Movie("Матрица", "Фантастика", 1999));
    tx.commit();
}
// ПОИСК ПО КЛЮЧУ И ЗАПРОС НА HQL: имена класса и полей,
// а не имена таблицы и столбцов
try (Session session = factory.openSession()) {
    Movie m = session.get(Movie.class, 1);
    List<Movie> found = session.createQuery(
            "FROM Movie WHERE genre = :g ORDER BY year",
            Movie.class)
        .setParameter("g", "Фантастика").list();
}
factory.close();
\end{minted}
\caption{Операции Hibernate: сохранение, поиск по ключу, запрос на HQL}
\label{lst:06-9}
\end{listing}

\term{HQL} (\eng{Hibernate Query Language}) выглядит как SQL, но оперирует
именами классов и полей Java: в запросе стоит \texttt{FROM Movie}, а не
\texttt{FROM movies}; запросы на изменение выполняются отдельным методом
\texttt{createMutationQuery}. Тот же запрос можно собрать и из объектов
Java~--- через \term{Criteria API}, где вместо строки используются
\texttt{CriteriaBuilder}, \texttt{CriteriaQuery} и \texttt{Root}. Строк нет,
значит, опечатку поймает компилятор, а условия можно добавлять в цикле~--- это
главное преимущество при построении динамических фильтров. Расплата~---
многословность: запрос на HQL занимает строку, а на Criteria~API~--- пять.

Связи между таблицами описывают четыре аннотации: \texttt{@OneToOne}
(фильм~--- описание), \texttt{@ManyToOne} (фильм~--- режиссёр),
\texttt{@OneToMany} (режиссёр~--- фильмы) и \texttt{@ManyToMany} (фильмы~---
актёры). Первые две ставят вместе с \texttt{@JoinColumn} на владеющей стороне,
в таблице которой лежит внешний ключ; \texttt{@OneToMany} описывает ту же
связь с обратной стороны и обязана нести параметр \texttt{mappedBy} с именем
поля владельца, иначе в схеме появится лишний внешний ключ; для
\texttt{@ManyToMany} нужна отдельная таблица-связка, задаваемая аннотацией
\texttt{@JoinTable}. Параметр \texttt{fetch = FetchType.LAZY} включает
\term{ленивую загрузку}: связанный объект не читается из базы, пока к полю не
обратятся, зато после закрытия сессии такое обращение даёт
\texttt{LazyInitializationException}.

\subsection{JPA и Hibernate, персистентный контекст}

Аннотации мы импортируем из пакета \texttt{jakarta.persistence}, а
\texttt{Session}~--- из \texttt{org.hibernate}: перед нами два уровня.
\term{JPA} (\eng{Jakarta Persistence API})~--- спецификация: набор
интерфейсов, аннотаций и правил, в котором нет кода, ходящего в базу;
Hibernate~--- её реализация. Разница такая же, как между строительными нормами
и бригадой, которая по ним строит: пока вы пользуетесь стандартными
аннотациями, бригаду можно сменить. У каждого понятия есть двойник: паре
\texttt{EntityManagerFactory} и \texttt{EntityManager} соответствуют
\texttt{SessionFactory} и \texttt{Session}, языку JPQL~--- HQL, файлу
\texttt{persistence.xml}~--- \texttt{hibernate.cfg.xml}.

Сессия~--- не просто трубка к базе: внутри неё живёт \term{персистентный
контекст} (\eng{persistence context}), набор объектов, за которыми она следит.
Это кэш первого уровня~--- запросив в одной сессии сущность с одним ключом
дважды, вы получите тот же объект. Аналогия~--- черновик на столе: пока вы его
правите, в архив ничего не уходит, а правки уезжают одной пачкой при
\term{сбросе} (\eng{flush}), совпадающем с коммитом. Отсюда четыре состояния
сущности (табл.~\ref{tab:06-3}).

\begin{table}[htbp]
\caption{Четыре состояния сущности}
\label{tab:06-3}
\centering
\begin{tabular}{L{3.4cm}L{11.2cm}}
\toprule
Состояние & Что оно означает \\
\midrule
\eng{transient} & обычный объект после \texttt{new}: контекст о нём не знает, строки в базе нет \\
\eng{persistent} & объект под управлением контекста: изменение полей само превратится в \texttt{UPDATE} \\
\eng{detached} & был управляемым, но сессия закрылась: строка есть, за полями никто не следит \\
\eng{removed} & помечен к удалению: \texttt{DELETE} уйдёт при ближайшем сбросе \\
\bottomrule
\end{tabular}
\end{table}

Переходы выполняют методы сессии: \texttt{persist} делает объект управляемым,
\texttt{get} загружает строку и сразу берёт объект под управление,
\texttt{remove} помечает к удалению, закрытие сессии переводит всё в
\eng{detached}. Отсюда классическая ошибка новичка: правка объекта после
\texttt{session.close()} в базу не попадает. Спасает новая сессия и метод
\texttt{merge}, но с тонкостью: он не превращает переданный объект в
управляемый, а возвращает \emph{новый} объект-копию.

\subsection{Lombok: меньше шаблонного кода}

Вернёмся к листингу~\ref{lst:06-7}: там, где написано «геттеры, сеттеры и
\texttt{toString()} опущены», должно быть около десяти методов, ни в одном из
которых нет логики. Представьте нотариуса, который на каждой странице толстой
стопки ставит одну и ту же подпись: можно расписаться сто раз, а можно
вырезать печать. \term{Lombok}~--- такая печать для Java-класса: библиотека
работает как обработчик аннотаций и дописывает нужные методы прямо в файл
\texttt{.class}. Её аннотации импортируются из пакета \texttt{lombok}
(листинг~\ref{lst:06-10}).

\begin{listing}[H]
\begin{minted}{java}
@Entity
@Table(name = "movies")
@Getter                 // геттеры для всех полей
@Setter                 // сеттеры для всех полей
@NoArgsConstructor      // обязателен для JPA
@AllArgsConstructor     // конструктор со всеми полями
@ToString
public class Movie {
    // поля те же, что в листинге 6.7;
    // ни одного метода писать не нужно
}
\end{minted}
\caption{Класс-сущность с аннотациями Lombok}
\label{lst:06-10}
\end{listing}

Подключается Lombok зависимостью с областью видимости \texttt{provided}: при
компиляции нужна, в готовый JAR не попадает. Пятьдесят строк сжались до
пятнадцати, а поведение не изменилось.

\begin{vrezka}
\textbf{Важно.} Конструктор без аргументов обязателен для JPA, а не просто
удобен: Hibernate создаёт объект сущности через рефлексию \emph{до} того, как
заполнит поля значениями из базы. А аннотацию \texttt{@Data} на сущность
вешать не стоит: её \texttt{equals} и \texttt{toString} используют все поля,
включая ленивые связи, и вывод объекта в консоль обернётся десятком лишних
запросов или исключением.
\end{vrezka}

\section{Транзакции и уровни изоляции}

\subsection{Свойства ACID и управление транзакцией}

\term{Транзакция}~--- группа операций с базой, выполняемая как единое
неделимое целое. Представьте покупку двух билетов в кино на соседние места:
либо вам достаются оба, либо ни одного. Требования к транзакциям сокращают в
аббревиатуру \term{ACID}: \emph{атомарность}~--- всё или ничего, если на
середине что-то упало, база откатывает уже сделанное;
\emph{согласованность}~--- ограничения целостности не нарушаются;
\emph{изолированность}~--- параллельные транзакции не мешают друг другу;
\emph{долговечность}~--- после фиксации данные переживут отключение питания.
Три свойства из четырёх обеспечиваются сами; изолированность~--- регулятор,
который вы поворачиваете сами, выбирая между строгостью и скоростью.

По умолчанию соединение работает в режиме автоматической фиксации: каждый
запрос сам себе транзакция. Чтобы объединить запросы, режим выключают вызовом
\texttt{setAutoCommit(false)}, а дальше сами решают, завершить транзакцию
методом \texttt{commit} или отменить методом \texttt{rollback}. Иногда
отменить хочется только последний кусок работы~--- для этого есть \term{точка
сохранения} (\eng{savepoint}), работающая как сохранение в компьютерной игре
(листинг~\ref{lst:06-11}).

\begin{listing}[H]
\begin{minted}{java}
// Заказ обязателен, начисление бонусов - необязательно
conn.setAutoCommit(false);          // начало транзакции
try {
    insertOrderRow(conn, customerId, productId);

    // До сюда заказ уже корректен - ставим отметку
    Savepoint afterOrder = conn.setSavepoint("after_order");
    try {
        addBonusPoints(conn, customerId, bonus);
    } catch (SQLException bonusError) {
        conn.rollback(afterOrder);  // откат до отметки
    }
    conn.commit();                  // фиксируем то, что уцелело
} catch (SQLException e) {
    conn.rollback();                // полный откат транзакции
    throw e;
} finally {
    conn.setAutoCommit(true);       // вернуть режим как было
}
\end{minted}
\caption{Транзакция с точкой сохранения}
\label{lst:06-11}
\end{listing}

Блок \texttt{finally} здесь не формальность: соединение почти всегда приходит
из \term{пула соединений} (\eng{connection pool})~--- набора заранее открытых
соединений, которые выдаются на время и возвращаются обратно, как книги в
библиотеке. Оставить возвращённое соединение в режиме открытой транзакции
значит подложить мину следующему, кто его возьмёт.

\subsection{Аномалии параллельного выполнения}

Пока транзакции идут по очереди, всё безупречно. Проблемы начинаются, когда
они пересекаются во времени. Классических аномалий четыре
(табл.~\ref{tab:06-4}); представляйте себе две транзакции, работающие с одним
и тем же счётом.

\begin{table}[htbp]
\caption{Аномалии параллельного выполнения транзакций}
\label{tab:06-4}
\centering
\begin{tabular}{L{3.5cm}L{7.3cm}L{4cm}}
\toprule
Аномалия & Что происходит & Запрещена с уровня \\
\midrule
Грязное чтение & транзакция видит ещё не зафиксированные изменения соседней; после её отката окажется, что таких данных не было & \eng{READ COMMITTED} \\
Неповторяющееся чтение & два одинаковых чтения одной строки внутри транзакции дают разные значения & \eng{REPEATABLE READ} \\
Фантомное чтение & между двумя одинаковыми выборками меняется состав строк: появляются новые & \eng{SERIALIZABLE} \\
Потерянное обновление & две транзакции прочитали одно значение и записали своё; вторая затёрла работу первой & устраняется не уровнем, а блокировкой \\
\bottomrule
\end{tabular}
\end{table}

Разница между вторым и третьим случаем тонкая: при неповторяющемся чтении
меняется содержимое известной строки, при фантомном~--- сам состав выборки.
Потерянное обновление стоит особняком: от него можно защититься, не трогая
уровень изоляции,~--- достаточно писать
\texttt{UPDATE accounts SET balance = balance + 500} вместо чтения значения в
программу и обратной записи.

\subsection{Четыре уровня изоляции}

Стандарт SQL определяет четыре уровня, и каждый следующий запрещает больше
аномалий, но стоит дороже: чем строже изоляция, тем больше блокировок и тем
чаще транзакции ждут друг друга (табл.~\ref{tab:06-5}).

\begin{table}[htbp]
\caption{Уровни изоляции и допускаемые ими аномалии}
\label{tab:06-5}
\centering
\begin{tabular}{L{4.6cm}C{3.1cm}C{3.5cm}C{2.8cm}}
\toprule
Уровень & Грязное чтение & Неповторяющееся чтение & Фантомное чтение \\
\midrule
\eng{READ UNCOMMITTED} & возможно & возможно & возможно \\
\eng{READ COMMITTED} & нет & возможно & возможно \\
\eng{REPEATABLE READ} & нет & нет & возможно \\
\eng{SERIALIZABLE} & нет & нет & нет \\
\bottomrule
\end{tabular}
\end{table}

Запомнить их проще всего через бытовую сцену: вы читаете документ, над которым
работает коллега. Заглядывать через плечо во время набора~---
\eng{READ UNCOMMITTED}: быстро, но вы читаете и то, что он сотрёт. Смотреть
только сохранённые версии~--- \eng{READ COMMITTED}. Взять копию и читать
только её~--- \eng{REPEATABLE READ}. Запретить коллеге трогать документ~---
\eng{SERIALIZABLE}: надёжнее всего, но коллега простаивает. Уровень задаётся
на соединении и обязательно \emph{до} начала транзакции
(листинг~\ref{lst:06-12}).

\begin{listing}[H]
\begin{minted}{java}
// Уровень ставим ДО setAutoCommit и до первого запроса
conn.setTransactionIsolation(
        Connection.TRANSACTION_REPEATABLE_READ);
conn.setAutoCommit(false);
\end{minted}
\caption{Установка уровня изоляции в JDBC}
\label{lst:06-12}
\end{listing}

Остальные константы того же интерфейса~---
\texttt{TRANSACTION\_READ\_UNCOMMITTED}, \texttt{TRANSACTION\_READ\_COMMITTED}
и \texttt{TRANSACTION\_SERIALIZABLE}; узнать уровень по умолчанию позволяет
\texttt{conn.getMetaData().getDefaultTransactionIsolation()}. Тонкость в том,
что JDBC описывает интерфейс, а поддерживает его конкретная СУБД: PostgreSQL
молча выполнит \eng{READ UNCOMMITTED} как \eng{READ COMMITTED}, MySQL по
умолчанию работает на уровне \eng{REPEATABLE READ}, а Oracle реализует лишь
два уровня из четырёх. Не полагайтесь на умолчания: при переезде между СУБД
поведение приложения изменится молча.

\subsection{Блокировки в Hibernate}

Поднимать уровень изоляции ради защиты от потерянного обновления дорого;
дешевле \term{оптимистичная блокировка}: будем считать, что конфликтов почти
не бывает, но если конфликт случится~--- заметим его и не дадим затереть чужую
работу. Достаточно добавить в сущность поле с аннотацией \texttt{@Version}
(листинг~\ref{lst:06-13}).

\begin{listing}[H]
\begin{minted}{java}
@Entity
public class Account {
    @Id
    private Integer id;
    private long balance;

    @Version           // служебное поле, руками не трогаем
    private int version;
}

// Hibernate добавит версию в условие обновления:
// UPDATE accounts SET balance = ?, version = 6
//   WHERE id = ? AND version = 5
try {
    tx.commit();
} catch (StaleStateException | OptimisticLockException e) {
    System.out.println("Данные изменил кто-то другой");
}
\end{minted}
\caption{Оптимистичная блокировка через версию строки}
\label{lst:06-13}
\end{listing}

Если параллельная транзакция уже зафиксировала свои изменения, версия в базе
стала равна шести, условие \texttt{version = 5} не выполнится, обновится ноль
строк~--- и коммит завершится исключением. Это как бронирование места в
поезде: вы спокойно выбираете вагон и полку, и только при оплате система
проверяет, свободно ли ещё выбранное место.

Если конфликты не редкость, а норма~--- склад с остатками товара, продажа
последних билетов,~--- выгоднее \term{пессимистичная блокировка}: строка
захватывается сразу при чтении вызовом
\texttt{session.get(Account.class, 1, LockMode.PESSIMISTIC\_WRITE)}, что
превращается в запрос \texttt{SELECT ... FOR UPDATE}, и остальные транзакции
ждут до конца вашей~--- за это платят простоем и риском взаимной блокировки.
Всё, что здесь написано руками, в Spring сводится к аннотации
\texttt{@Transactional} над методом сервиса; она разбирается в
главе~\ref{ch:07}.

\section{Функциональный стиль: Stream API}

\subsection{Поток данных против коллекции}

Коллекция~--- это ящик с яблоками в кладовке: она отвечает на вопросы «что у
меня есть» и «сколько». \term{Поток} (\eng{stream})~--- конвейерная лента на
сортировочной линии: сами яблоки на ней не живут, они по ней едут, а вдоль
ленты стоят рабочие~--- один выбрасывает подгнившие, второй наклеивает
стикеры, третий раскладывает по ящикам. Лента ничего не хранит: она описывает,
что делать с каждым яблоком. Отсюда все отличия: коллекция хранит элементы,
поток~--- нет; коллекцию можно обойти сколько угодно раз, поток~--- ровно
один; поток ленив и не делает ничего, пока не потребуют результат
(листинг~\ref{lst:06-14}).

\begin{listing}[H]
\begin{minted}{java}
// Императивно: вы командуете компьютером пошагово
List<String> namesA = new ArrayList<>();
for (User user : users) {
    if (user.age() >= 18) {
        namesA.add(user.name());
    }
}
Collections.sort(namesA);

// Декларативно: вы описываете нужный результат
List<String> namesB = users.stream()
        .filter(user -> user.age() >= 18)
        .map(User::name)
        .sorted()
        .toList();
\end{minted}
\caption{Императивный и декларативный стиль}
\label{lst:06-14}
\end{listing}

Второй вариант короче, но короче~--- не главное: в нём нет
переменной-накопителя, которую можно испортить, нет индексов, в которых можно
ошибиться, и каждая строка читается как отдельная фраза. Любой конвейер
состоит ровно из трёх частей (рис.~\ref{fig:06-3}).

\begin{figure}[htbp]\centering
\begin{tikzpicture}[font=\small,node distance=6mm]
\node[boxw=32mm] (a) {\textbf{Источник}\\[2pt]\scriptsize коллекция, массив,\\\scriptsize файл, генератор};
\node[boxw=42mm,right=of a] (b) {\textbf{Промежуточные}\\[2pt]\scriptsize \texttt{filter}, \texttt{map}, \texttt{sorted}\\\scriptsize ленивы, возвращают поток};
\node[boxw=42mm,right=of b] (c) {\textbf{Терминальная}\\[2pt]\scriptsize \texttt{collect}, \texttt{count}\\\scriptsize запускает весь конвейер};
\draw[flowarrow] (a)--(b);\draw[flowarrow] (b)--(c);
\end{tikzpicture}
\caption{Строение конвейера обработки данных}\label{fig:06-3}
\end{figure}

Не путайте термины: в главе~\ref{ch:05} словом «поток» назывались поток
выполнения (\texttt{Thread}) и поток ввода-вывода (\texttt{InputStream}).
\texttt{Stream} из пакета \texttt{java.util.stream}~--- третья, независимая
сущность.

\subsection{Создание потока и промежуточные операции}

Источником потока может быть коллекция, массив, перечисленные значения,
правило генерации и даже файл (листинг~\ref{lst:06-15}). Массив превращает в
поток метод \texttt{Arrays.stream}. Отдельно стоит последний случай:
\texttt{Files.lines} читает файл по строке по мере надобности, поэтому за
потоком стоит открытый файловый дескриптор и закрывать такой поток
обязательно.

\begin{listing}[H]
\begin{minted}{java}
// Из коллекции и из перечисленных значений
List<String> cities = List.of("Москва", "Казань", "Сочи");
Stream<String> s1 = cities.stream();
Stream<String> s2 = Stream.of("а", "б", "в");

// Бесконечный поток по правилу: без limit он не кончится
List<Integer> powers = Stream.iterate(1, n -> n * 2)
        .limit(10).toList();          // [1, 2, 4, ..., 512]

// Строки файла читаются лениво, поэтому поток закрывают
try (Stream<String> lines = Files.lines(Path.of("app.log"))) {
    long errors = lines.filter(l -> l.contains("ERROR")).count();
    System.out.println("Строк с ошибками: " + errors);
}
\end{minted}
\caption{Способы получить поток}
\label{lst:06-15}
\end{listing}

Промежуточная операция возвращает новый поток, поэтому такие операции
выстраиваются в цепочку любой длины~--- каждая из них очередной рабочий на
конвейере (табл.~\ref{tab:06-6}).

\begin{table}[htbp]
\caption{Промежуточные операции потока}
\label{tab:06-6}
\centering
\begin{tabular}{L{4.4cm}L{10.2cm}}
\toprule
Операция & Что делает \\
\midrule
\texttt{filter(p)} & пропускает дальше элементы, для которых предикат истинен \\
\texttt{map(f)} & преобразует каждый элемент в другой, возможно другого типа \\
\texttt{flatMap(f)} & превращает каждый элемент в поток и сливает их в один \\
\texttt{distinct()} & убирает дубликаты, сравнивая через \texttt{equals} \\
\texttt{sorted(cmp)} & сортирует по естественному порядку или по компаратору \\
\texttt{limit(n)}, \texttt{skip(n)} & оставляет первые \texttt{n} элементов, пропускает первые \texttt{n} \\
\texttt{takeWhile(p)} & берёт элементы, пока предикат истинен, и останавливается \\
\bottomrule
\end{tabular}
\end{table}

Типичная цепочка показана в листинге~\ref{lst:06-16}. Обратите внимание на две
последние строки: \texttt{takeWhile} обрывает поток на первом несовпадении, а
\texttt{filter} проверяет все элементы до конца.

\begin{listing}[H]
\begin{minted}{java}
record User(String name, int age, String city) {}

List<User> users = List.of(
        new User("Анна", 28, "Москва"),
        new User("Борис", 17, "Казань"),
        new User("Вера", 34, "Москва"),
        new User("Анна", 28, "Москва"));   // полный дубликат

List<String> names = users.stream()
        .filter(u -> u.age() >= 18)
        .distinct()            // record сам умеет equals
        .sorted(Comparator.comparingInt(User::age).reversed())
        .map(User::name)
        .toList();             // [Вера, Анна]

List<Integer> nums = List.of(1, 3, 5, 8, 7, 9);
nums.stream().takeWhile(n -> n % 2 != 0).toList(); // [1, 3, 5]
nums.stream().filter(n -> n % 2 != 0).toList();    // [1,3,5,7,9]
\end{minted}
\caption{Цепочки промежуточных операций}
\label{lst:06-16}
\end{listing}

Отдельного разговора заслуживает \texttt{flatMap}: он нужен, когда каждый
элемент сам является набором, а вам нужен один плоский список. Список списков
имён превращается в список имён вызовом
\texttt{teams.stream().flatMap(List::stream).toList()}. Разница такая:
\texttt{map}~--- «один элемент на входе, один на выходе»,
\texttt{flatMap}~--- «один на входе, ноль или больше на выходе».

\subsection{Ленивость~--- главное свойство потоков}

Вот утверждение, которое чаще всего понимают неправильно: промежуточные
операции не выполняются в момент вызова. Вызов \texttt{filter} не фильтрует,
вызов \texttt{map} не преобразует~--- они лишь достраивают описание конвейера.
Работа начинается, только когда вызвана терминальная операция: она включает
мотор и тянет элементы через всю цепочку. Диктуя официанту заказ, вы не ждёте,
что повар начнёт жарить после каждого вашего слова. Убедиться в этом позволяет
промежуточная операция \texttt{peek}: она подсматривает за проходящими
элементами, не меняя их (листинг~\ref{lst:06-17}).

\begin{listing}[H]
\begin{minted}{java}
List<String> cities = List.of("Москва", "Казань", "Самара");

Stream<String> pipeline = cities.stream()
        .peek(c -> System.out.println("  peek: " + c))
        .filter(c -> c.length() > 5)
        .map(String::toUpperCase);

System.out.println("Конвейер построен, но не запущен");
System.out.println(pipeline.toList());

// Вывод: сначала строка "Конвейер построен, но не запущен",
// и только затем три строки "peek: ..." и результат
// [МОСКВА, КАЗАНЬ, САМАРА]
\end{minted}
\caption{Ленивость конвейера, показанная через peek}
\label{lst:06-17}
\end{listing}

Из ленивости следует, что элементы обрабатываются «по вертикали»: поток берёт
первый элемент, проводит его через всю цепочку до конца, затем берётся за
второй. Поэтому конвейер с \texttt{findFirst} останавливается, как только
ответ найден. Операции \texttt{findFirst}, \texttt{findAny},
\texttt{anyMatch}, \texttt{allMatch}, \texttt{noneMatch} и \texttt{limit} за
это называют \term{короткозамыкающими} (\eng{short-circuiting}).

\begin{oshibka}
\textbf{Частая ошибка: конвейер без терминальной операции.} Строка
\texttt{users.stream().map(u -> u.name().toUpperCase());} ничего не изменит и
ничего не напечатает~--- она лишь построит описание и выбросит его. И не
пишите в \texttt{peek} ничего, кроме печати для отладки: спецификация не
гарантирует, что он будет вызван для каждого элемента.
\end{oshibka}

\subsection{Терминальные операции и сборка результата}

Терминальная операция завершает конвейер: она возвращает не поток, а
конкретный результат~--- или не возвращает ничего. Их немного:
\texttt{forEach} и \texttt{count} выполняют действие и считают элементы,
\texttt{min} и \texttt{max} находят крайние, \texttt{anyMatch},
\texttt{allMatch} и \texttt{noneMatch} отвечают на вопросы «хоть один», «все»,
«ни один», \texttt{findFirst} возвращает первый элемент, \texttt{reduce}
сворачивает поток в одно значение, а \texttt{toList} и \texttt{collect}
собирают результат в коллекцию (листинг~\ref{lst:06-18}).

\begin{listing}[H]
\begin{minted}{java}
List<Integer> numbers = List.of(4, 8, 15, 16, 23, 42);
int sum = numbers.stream().reduce(0, Integer::sum);

long adults = users.stream().filter(u -> u.age() >= 18).count();

Optional<User> oldest = users.stream()
        .max(Comparator.comparingInt(User::age));
oldest.ifPresent(u -> System.out.println(u.name()));
\end{minted}
\caption{Свёртка, подсчёты и поиск крайнего}
\label{lst:06-18}
\end{listing}

Одна ловушка стоит того, чтобы узнать о ней заранее: на \emph{пустом} потоке
\texttt{allMatch} и \texttt{noneMatch} возвращают \texttt{true}, а
\texttt{anyMatch}~--- \texttt{false}. Это не ошибка, а математическая логика:
утверждение «все элементы пустого множества обладают свойством X» истинно,
потому что опровергнуть его нечем. Поэтому проверку \texttt{allMatch} почти
всегда сопровождают проверкой, что данные вообще есть.

Самая мощная терминальная операция~--- \texttt{collect}: в неё передаётся
\term{сборщик} (\eng{collector}), описывающий, как накапливать результат.
Готовые сборщики живут в классе \texttt{Collectors}, и самый интересный из
них~--- \texttt{groupingBy}: он раскладывает элементы по корзинам, как
почтальон раскладывает письма по ячейкам подъезда. Первый аргумент определяет
ячейку, а второй~--- \term{нижестоящий сборщик} (\eng{downstream collector})~---
говорит, что сделать с письмами внутри каждой ячейки
(листинг~\ref{lst:06-19}).

\begin{listing}[H]
\begin{minted}{java}
record Employee(String name, String dept, int salary) {}

List<Employee> staff = List.of(
    new Employee("Анна",  "Разработка",   180_000),
    new Employee("Борис", "Разработка",   210_000),
    new Employee("Вера",  "Тестирование", 140_000));

// Сколько человек в каждом отделе
Map<String, Long> byDept = staff.stream()
        .collect(Collectors.groupingBy(Employee::dept,
                 Collectors.counting()));

// Средняя зарплата по отделам: другой нижестоящий сборщик
Map<String, Double> avg = staff.stream()
        .collect(Collectors.groupingBy(Employee::dept,
                 Collectors.averagingInt(Employee::salary)));

// В корзины кладём не сотрудников, а только их имена;
// третий аргумент TreeMap::new отсортировал бы ключи
Map<String, List<String>> byName = staff.stream()
        .collect(Collectors.groupingBy(Employee::dept,
                 Collectors.mapping(Employee::name,
                                    Collectors.toList())));

// partitioningBy - ровно две корзины по предикату
Map<Boolean, List<Employee>> rich = staff.stream()
        .collect(Collectors.partitioningBy(
                 e -> e.salary() >= 160_000));
\end{minted}
\caption{Сборка, группировка и разбиение}
\label{lst:06-19}
\end{listing}

Про печать результатов договоримся сразу: \texttt{groupingBy} собирает
результат в \texttt{HashMap}, а она порядок ключей не хранит~--- пары почти
наверняка окажутся не в том порядке, в каком шли элементы; нужен порядок~---
задайте тип карты. Последний в листинге \texttt{partitioningBy}~--- частный
случай группировки ровно на две корзины по предикату; его преимущество в том,
что в результате всегда есть оба ключа, даже если одна корзина пуста. Ещё один
частый сборщик, \texttt{Collectors.joining(", ", "[", "]")}, склеивает поток
строк в одну с разделителем, префиксом и суффиксом.

\subsection{Примитивные потоки и Optional}

\texttt{Stream<Integer>} хранит объекты, а значит, на каждое число тратится
упаковка в \texttt{Integer}; на миллионе элементов это заметно. Поэтому есть
три специализированных потока~--- \texttt{IntStream}, \texttt{LongStream} и
\texttt{DoubleStream}. Переход в них делают методы \texttt{mapToInt} и его
родственники, обратно~--- метод \texttt{boxed}. Взамен они дают операции,
которых у объектного потока нет: \texttt{sum}, \texttt{average} и особенно
\texttt{summaryStatistics}, считающий за один проход сразу пять величин
(листинг~\ref{lst:06-20}).

\begin{listing}[H]
\begin{minted}{java}
int total = staff.stream().mapToInt(Employee::salary).sum();

// range даёт 0..4, rangeClosed - 1..5; boxed возвращает
// обратно в Stream<Integer>
List<Integer> boxed = IntStream.rangeClosed(1, 5)
        .boxed().toList();

IntSummaryStatistics st = staff.stream()
        .mapToInt(Employee::salary)
        .summaryStatistics();
System.out.printf("От %d до %d, в среднем %.2f%n",
        st.getMin(), st.getMax(), st.getAverage());
\end{minted}
\caption{Примитивный поток и сводная статистика}
\label{lst:06-20}
\end{listing}

Заметьте: \texttt{average} возвращает не число, а \texttt{OptionalDouble}~---
у пустого потока нет ни среднего, ни максимума, и язык заставляет обработать
этот случай явно. \texttt{Optional<T>}~--- контейнер, который либо содержит
значение, либо пуст; это способ сказать «результата может не быть» на уровне
типа, а не через \texttt{null} и \texttt{NullPointerException}. Значение
достают методами \texttt{orElse} (запасной вариант), \texttt{orElseThrow}
(своё исключение с внятным текстом) и \texttt{ifPresentOrElse} (два действия
на оба случая); метод \texttt{get} без проверки использовать не стоит.

\subsection{Параллельные потоки и одноразовость}

Одна из причин появления Stream~API~--- возможность распараллелить обработку,
не написав ни строчки про потоки выполнения: достаточно заменить
\texttt{stream()} на \texttt{parallelStream()}. Поток разбивается на части,
части считаются в общем пуле \texttt{ForkJoinPool}, результаты сливаются
обратно. Бесплатным это не бывает: параллельность помогает, когда данных
много, источник легко делится (массив, \texttt{ArrayList}, диапазон чисел), а
операции тяжёлые и независимые; и вредит, когда коллекция мала, источник
делится плохо (\texttt{LinkedList}, файл) или важен порядок. Главный
запрет~--- никакого разделяемого изменяемого состояния
(листинг~\ref{lst:06-21}).

\begin{listing}[H]
\begin{minted}{java}
// ОШИБКА: ArrayList не потокобезопасен. Результат -
// потерянные элементы или исключение при записи
List<Integer> unsafe = new ArrayList<>();
IntStream.range(0, 100_000).parallel().forEach(unsafe::add);

// ПРАВИЛЬНО: пусть поток соберёт результат сам
List<Integer> safe = IntStream.range(0, 100_000)
        .parallel().boxed().toList();

// ОШИБКА: поток одноразовый. Второй вызов бросит
// IllegalStateException; переиспользуют источник, а не поток
Stream<String> stream = cities.stream();
stream.forEach(System.out::println);
stream.forEach(System.out::println);
\end{minted}
\caption{Ошибки параллельного и повторного использования потока}
\label{lst:06-21}
\end{listing}

Одноразовость~--- прямое следствие природы потока: он ничего не хранит и лишь
тянет элементы из источника; пройти по нему второй раз~--- всё равно что
попросить конвейерную ленту прокрутить назад уже уехавшие яблоки. И совет
напоследок: не ставьте \texttt{parallel()} «на всякий случай»~--- сначала
напишите последовательную версию и измерьте время.

\praktikum{}

Понадобится JDK~17 или новее, установленный Maven и подключение к сети при
первой сборке. База данных отдельной установки не требует: H2 подключается как
обычная библиотека.

\subsection*{Задание 1. Maven-проект и его зависимости}

Создайте вручную, без помощи среды разработки, каталог \texttt{movie-app} со
структурой из рис.~\ref{fig:06-1} и файлом \texttt{pom.xml} по образцу
листинга~\ref{lst:06-1}, положите в \texttt{src/main/java/com/movies} класс
\texttt{Main} и выполните по очереди команды из листинга~\ref{lst:06-2},
записывая, что изменилось в каталоге \texttt{target} после каждой. Затем
добавьте зависимость от \texttt{hibernate-core} версии 6.4.0.Final, выполните
\texttt{mvn dependency:tree} и ответьте письменно, сколько у проекта прямых
зависимостей и сколько транзитивных добавил Hibernate.

\subsection*{Задание 2. CRUD через JDBC и паттерн DAO}

Наберите программу из листинга~\ref{lst:06-4}, дополнив её вставкой четырёх
фильмов, обновлением одного из них, удалением ещё одного и выводом всех
оставшихся строк. Добавьте поиск по части названия запросом вида
\texttt{SELECT * FROM movies WHERE LOWER(title) LIKE LOWER(?)} и объясните
письменно, зачем в параметре нужны знаки процента по краям.

Вторая часть~--- перенос SQL в отдельный слой. По образцу
листинга~\ref{lst:06-6} напишите класс \texttt{Movie} (поля \texttt{id},
\texttt{title}, \texttt{genre}, \texttt{year}), интерфейс \texttt{MovieDAO} с
методами \texttt{insert}, \texttt{findById}, \texttt{findAll},
\texttt{findByGenre}, \texttt{update}, \texttt{delete} и его реализацию на
JDBC. Все методы обязаны использовать \texttt{PreparedStatement},
\texttt{insert}~--- забирать сгенерированный ключ, а преобразование строки
результата в объект должно быть вынесено в единственный приватный метод.
Напишите класс-проверку, который подключается к базе
\texttt{jdbc:h2:mem:movietest;DB\_CLOSE\_DELAY=-1}, вставляет четыре фильма и
вызывает каждый метод DAO хотя бы один раз.

\subsection*{Задание 3. Hibernate и Lombok}

Создайте файл \texttt{src/main/resources/hibernate.cfg.xml} по
листингу~\ref{lst:06-8} и перепишите класс \texttt{Movie} с аннотациями JPA по
листингу~\ref{lst:06-7}. Воспроизведите операции из листинга~\ref{lst:06-9}:
сохраните пять фильмов, найдите один по ключу, обновите год выхода запросом на
HQL через \texttt{createMutationQuery}, удалите фильм и выведите все фильмы
одного жанра; добавьте постраничный вывод методами \texttt{setFirstResult} и
\texttt{setMaxResults}. Оставьте включённым \texttt{hibernate.show\_sql} и
выпишите из консоли запросы, которые Hibernate сгенерировал за вас. Затем
подключите Lombok зависимостью с областью видимости \texttt{provided} и
перепишите \texttt{Movie} по листингу~\ref{lst:06-10}, не меняя ни строчки в
использующей его программе.

Ответьте письменно: чем HQL отличается от SQL; что произойдёт, если сменить
\texttt{hbm2ddl.auto} со значения \texttt{create} на \texttt{validate};
сколько строк исчезло из класса после подключения Lombok и почему конструктор
без аргументов для JPA обязателен.

\subsection*{Задание 4. Транзакции и уровни изоляции}

Создайте таблицу
\texttt{accounts (id INT PRIMARY KEY, owner VARCHAR(100), balance DECIMAL(10,2))}
и заполните её несколькими счетами. Напишите метод перевода средств, который
выключает автоматическую фиксацию, проверяет достаточность средств, списывает
и зачисляет сумму, фиксирует транзакцию при успехе, откатывает её при ошибке
или нехватке денег и восстанавливает исходный режим в блоке \texttt{finally}
по образцу листинга~\ref{lst:06-11}. Убедитесь, что после неудачного перевода
общая сумма на счетах не изменилась.

Вторая часть~--- эксперимент с изоляцией. Откройте к одной базе два соединения
и воспроизведите две аномалии из табл.~\ref{tab:06-4}. Для грязного чтения
поставьте на втором соединении \texttt{TRANSACTION\_READ\_UNCOMMITTED}, в
первом начните транзакцию, измените баланс и не фиксируйте её; прочитайте
баланс вторым соединением, затем откатите первое и прочитайте снова. Для
неповторяющегося чтения прочитайте баланс первым соединением, измените и
зафиксируйте его вторым, прочитайте первым повторно~--- и повторите тот же
сценарий на уровнях \texttt{TRANSACTION\_READ\_COMMITTED} и
\texttt{TRANSACTION\_REPEATABLE\_READ}. Уровень задавайте до вызова
\texttt{setAutoCommit(false)}, как в листинге~\ref{lst:06-12}. Ответьте
письменно, какие значения вернули оба чтения в каждом сценарии и почему они
различаются.

\subsection*{Задание 5. Обработка данных через Stream API}

Создайте класс \texttt{FilmStreams} с записью и списком фильмов из
листинга~\ref{lst:06-22} и решите на нём восемь задач, каждую~--- одним
конвейером, без единого цикла \texttt{for}.

\begin{listing}[H]
\begin{minted}{java}
record Film(String title, String genre, int year, double rating) {}

static final List<Film> FILMS = List.of(
    new Film("Матрица",       "Фантастика", 1999, 8.5),
    new Film("Начало",        "Фантастика", 2010, 8.7),
    new Film("Интерстеллар",  "Фантастика", 2014, 8.6),
    new Film("Тёмный рыцарь", "Боевик",     2008, 9.0),
    new Film("Джокер",        "Драма",      2019, 8.4),
    new Film("Паразиты",      "Триллер",    2019, 8.5),
    new Film("Зелёная миля",  "Драма",      1999, 9.1),
    new Film("Мементо",       "Триллер",    2000, 8.4));
\end{minted}
\caption{Исходные данные для заданий по Stream API}
\label{lst:06-22}
\end{listing}

Задачи такие: посчитать, сколько фильмов вышло до 2010~года; найти первый
фильм с рейтингом выше 9,0 и напечатать «не найден», если такого нет;
получить три самых свежих фильма; собрать все жанры без повторов по алфавиту;
посчитать средний рейтинг по жанрам; собрать названия фильмов по жанрам;
найти лучший фильм каждого жанра; разделить фильмы на вышедшие до 2010~года и
позже. Отдельно выведите сводную статистику по годам через
\texttt{summaryStatistics}. Затем свяжите обе половины занятия: возьмите
список, который возвращает \texttt{findAll} из задания~2 или запрос на HQL из
задания~3, и посчитайте по нему количество фильмов по жанрам и средний год
выхода с обработкой случая пустого списка. Ответьте письменно: что напечатает
конвейер, если убрать из него завершающий \texttt{toList()}; что вернёт
\texttt{allMatch} на пустом списке и почему; какой тип \texttt{Map} возвращает
\texttt{groupingBy} без третьего аргумента.

\subsection*{Результаты занятия}

По итогам занятия у вас должны быть готовы: проект \texttt{movie-app} с
рабочим \texttt{pom.xml} и отчётом по дереву зависимостей; программа CRUD на
чистом JDBC и комплект классов DAO; сущность с аннотациями JPA, файл
\texttt{hibernate.cfg.xml} и программа с запросами на HQL; вариант сущности на
Lombok; программа перевода средств и отчёт об эксперименте с уровнями
изоляции; класс \texttt{FilmStreams} с решениями всех восьми задач и
письменные ответы на вопросы заданий.

\kontrolvoprosy

\begin{enumerate}
\item Что такое GAV-координаты и зачем система сборки требует их у каждого
      проекта и каждой зависимости?
\item Чем область видимости \texttt{compile} отличается от \texttt{test} и
      \texttt{provided}? Приведите пример зависимости для каждой.
\item Что такое транзитивные зависимости и почему они могут стать проблемой?
\item Чем \texttt{PreparedStatement} отличается от \texttt{Statement} и
      почему подстановка пользовательского ввода в текст запроса недопустима?
\item Какие три выгоды даёт паттерн DAO? Почему метод поиска по ключу
      возвращает \texttt{Optional}, а не саму сущность?
\item Как соотносятся JPA и Hibernate? Чем \texttt{SessionFactory} отличается
      от \texttt{Session}?
\item Перечислите четыре состояния сущности. Почему изменение объекта после
      закрытия сессии не попадает в базу?
\item Что означают четыре буквы аббревиатуры ACID? Какое из свойств
      разработчик настраивает сам?
\item Назовите четыре аномалии параллельного выполнения. Какой уровень
      изоляции запрещает каждую из первых трёх и почему четвёртая стоит
      особняком?
\item В чём разница между оптимистичной и пессимистичной блокировкой?
\item Чем поток отличается от коллекции и почему по одному потоку нельзя
      пройти дважды?
\item Что такое ленивость промежуточных операций? Чем \texttt{map} отличается
      от \texttt{flatMap}?
\end{enumerate}

\testzadaniya

\begin{enumerate}

\vopros{Что такое GAV-координаты в Maven?}
  {\texttt{goal}, \texttt{action}, \texttt{value}~--- параметры плагина}
  {\texttt{groupId}, \texttt{artifactId}, \texttt{version}~--- адрес артефакта}
  {\texttt{git}, \texttt{archive}, \texttt{validate}~--- команды репозитория}
  {\texttt{gradle}, \texttt{ant}, \texttt{version}~--- совместимость систем}

\vopros{Что произойдёт при выполнении команды \texttt{mvn package}?}
  {выполнится только упаковка в JAR, без компиляции и тестов}
  {выполнятся \texttt{validate}, \texttt{compile}, \texttt{test} и \texttt{package}}
  {артефакт будет опубликован в центральном репозитории}
  {будут выполнены только \texttt{package} и \texttt{deploy}}

\vopros{Какая область видимости нужна библиотеке, которая используется только
        в тестах и не должна попасть в готовый JAR?}
  {\texttt{compile}}
  {\texttt{provided}}
  {\texttt{runtime}}
  {\texttt{test}}

\vopros{Чем \texttt{PreparedStatement} отличается от \texttt{Statement}?}
  {работает только с запросами \texttt{SELECT}}
  {не поддерживает параметры, зато выполняется быстрее}
  {выполняет запрос, не открывая соединения с базой}
  {отправляет шаблон и параметры отдельно, чем защищает от инъекции}

\vopros{Как соотносятся Hibernate и JPA?}
  {Hibernate~--- реализация спецификации JPA}
  {JPA~--- реализация Hibernate}
  {это два названия одного и того же продукта}
  {JPA пришёл на смену Hibernate и заменил его}

\vopros{Чем \texttt{SessionFactory} отличается от \texttt{Session}?}
  {оба объекта создаются заново на каждую операцию с базой}
  {\texttt{SessionFactory} создаётся на каждую операцию, \texttt{Session}~--- один раз}
  {разницы нет, это синонимы}
  {\texttt{SessionFactory} тяжёлый и создаётся один раз, \texttt{Session}~--- лёгкий, на единицу работы}

\vopros{Почему конструктор без аргументов обязателен для JPA-сущности?}
  {язык Java требует его от любого класса}
  {без него не компилируется аннотация \texttt{@Entity}}
  {без него сущность не попадёт в файл конфигурации}
  {Hibernate создаёт объект рефлексией до того, как заполнит поля}

\vopros{Что гарантирует свойство атомарности транзакции?}
  {операции выполнятся все либо ни одна}
  {параллельные транзакции не увидят работу друг друга}
  {данные переживут отключение питания}
  {запросы будут выполнены в порядке их отправки}

\vopros{Какой минимальный уровень изоляции гарантирует, что повторное чтение
        одной строки внутри транзакции вернёт то же значение?}
  {\eng{READ UNCOMMITTED}}
  {\eng{READ COMMITTED}}
  {\eng{REPEATABLE READ}}
  {\eng{SERIALIZABLE}}

\vopros{Что делает аннотация \texttt{@Version} в сущности Hibernate?}
  {хранит номер версии схемы базы для миграций}
  {включает оптимистичную блокировку: версия попадает в условие \texttt{UPDATE}}
  {включает пессимистичную блокировку строки при чтении}
  {задаёт версию Hibernate, с которой совместим класс}

\vopros{Что напечатает конвейер из \texttt{peek} и \texttt{filter}, у которого
        нет терминальной операции?}
  {ничего: без терминальной операции конвейер не запускается}
  {все элементы источника}
  {только первый элемент источника}
  {код не скомпилируется}

\vopros{Что окажется в списке после применения \texttt{flatMap(List::stream)} к
        потоку из списков \texttt{[Анна, Борис]} и \texttt{[Вера]}?}
  {\texttt{[Анна, Борис, Вера]}}
  {вложенность сохранится: \texttt{[[Анна, Борис], [Вера]]}}
  {по первому элементу каждого списка: \texttt{[Анна, Вера]}}
  {пустой список}

\vopros{Что вернёт \texttt{allMatch} на пустом потоке?}
  {\texttt{true}}
  {\texttt{false}}
  {исключение \texttt{NoSuchElementException}}
  {\texttt{Optional.empty()}}

\end{enumerate}
